\documentclass[a4paper,9pt]{extarticle}

% --- PACKAGES STANDARDS & COMPATIBILITÉ ATS MAXIMALE ---
\usepackage[utf8]{inputenc}
\usepackage[T1]{fontenc}
\usepackage{helvet}
\renewcommand{\familydefault}{\sfdefault}
\linespread{0.935}

% Marges équilibrées & sécurisées (0.42cm haut/bas, 1.1cm gauche/droite)
\usepackage[left=1.1cm, right=1.1cm, top=0.42cm, bottom=0.42cm]{geometry}
\usepackage{xcolor}
\usepackage{enumitem}
\usepackage{fontawesome5}
\usepackage{accsupp}
\usepackage[hidelinks]{hyperref}

\pagestyle{empty}

% --- PALETTE DE COULEURS MODERNE & ÉLÉGANTE ---
\definecolor{primary}{HTML}{0F172A}      % Slate 900 / Bleu marine profond (#0F172A)
\definecolor{accent}{HTML}{2563EB}       % Bleu royal électrique (#2563EB)
\definecolor{textdark}{HTML}{1E293B}     % Texte principal très lisible
\definecolor{textmuted}{HTML}{475569}    % Métadonnées / Dates / Lieux
\definecolor{rulecolor}{HTML}{CBD5E1}    % Lignes de séparation discrètes

\hypersetup{
    colorlinks=true,
    urlcolor=accent,
    linkcolor=accent
}

% --- MACRO ICÔNES SÉCURISÉES POUR LES ROBOTS ATS ---
% L'icône est affichée visuellement avec éclat mais est totalement masquée à l'extracteur de texte
\newcommand{\safeicon}[1]{\BeginAccSupp{ActualText={}}#1\EndAccSupp{}}

% Macro pour titres de section modernes avec icône & séparateur propre
\newcommand{\cvsection}[2]{%
  \vspace{1.5pt}%
  \noindent{\color{primary}\bfseries\small\safeicon{\textcolor{accent}{#1}}\enspace\MakeUppercase{#2}}\\[-3.5pt]%
  {\color{rulecolor}\rule{\linewidth}{0.55pt}}\par\vspace{0.8pt}%
}

% --- LISTES À PUCES STANDARDS & COMPACTES ---
\setlist[itemize]{
  noitemsep,
  topsep=0.4pt,
  parsep=0pt,
  partopsep=0pt,
  leftmargin=9.5pt,
  label={\color{accent}\small\textbullet}
}

\setlength{\parindent}{0pt}
\setlength{\parskip}{0pt}

\begin{document}

%=============================================================================
% EN-TÊTE MODERNE & COORDONNÉES DU CANDIDAT (100% SÉCURISÉES ATS)
%=============================================================================
\begin{center}
  {\LARGE \textbf{\color{primary} ANAS MOUQINE}}\\[1.5pt]
  {\small \textbf{\color{accent} ÉLÈVE-INGÉNIEUR EN 5\textsuperscript{ÈME} ANNÉE -- INTELLIGENCE ARTIFICIELLE (IA) \& DATA ENGINEERING}}\\[1.5pt]
  {\footnotesize \color{textmuted} Recherche de Stage PFE (Projet de Fin d'Études) de 4 à 6 mois \textbar\ Disponible Début 2026}\\[2.5pt]
  {\footnotesize \color{textdark}
    \safeicon{\textcolor{accent}{\faMapMarker*}}\ Casablanca, Maroc \quad\textbar\quad
    \safeicon{\textcolor{accent}{\faPhone*}}\ +212 672-104126 \quad\textbar\quad
    \safeicon{\textcolor{accent}{\faEnvelope}}\ \href{mailto:anasmouqine@gmail.com}{anasmouqine@gmail.com}
  }\\[1.5pt]
  {\footnotesize \color{textdark}
    \safeicon{\textcolor{accent}{\faLinkedin}}\ \href{https://www.linkedin.com/in/anas-mouqine-329207304/}{linkedin.com/in/anas-mouqine} \quad\textbar\quad
    \safeicon{\textcolor{accent}{\faGithub}}\ \href{https://github.com/anasmouquinee}{github.com/anasmouquinee} \quad\textbar\quad
    \safeicon{\textcolor{accent}{\faGlobe}}\ \href{https://anasmouquinee.github.io/port/}{anasmouquinee.github.io/port}
  }
\end{center}

\vspace{-3.5pt}

%=============================================================================
% 1. PROFIL PROFESSIONNEL
%=============================================================================
\cvsection{\faUser}{Profil Professionnel}
{\footnotesize \color{textdark}
Élève-Ingénieur en 5\textsuperscript{ème} année à l'École Marocaine des Sciences de l'Ingénieur (EMSI), spécialisé en \textbf{Intelligence Artificielle (IA)} et \textbf{Data Engineering}. Passionné par la conception de systèmes multi-agents autonomes, de pipelines Big Data Lakehouse (Apache Airflow, Kafka) et d'architectures distribuées. Fort d'expériences industrielles concrètes (\textbf{COSUMAR}, \textbf{SEWS-CABINET}), je recherche un stage PFE (Projet de Fin d'Études) de 4 à 6 mois pour concevoir et déployer des solutions de données et d'IA à fort impact opérationnel.
}

%=============================================================================
% 2. COMPÉTENCES TECHNIQUES (ACRONYMES DÉVELOPPÉS POUR RECHERCHE SÉMANTIQUE ATS)
%=============================================================================
\cvsection{\faCubes}{Compétences Techniques}
\begin{itemize}
  \item \textbf{Intelligence Artificielle (IA) \& Machine Learning :} Systèmes Multi-Agents (Multi-Agent Systems), Retrieval-Augmented Generation (RAG), Grands Modèles de Langage (LLMs / Large Language Models), LangChain, Prompt Engineering, Recherche Vectorielle (Vector Search), Apprentissage Automatique.
  \item \textbf{Big Data \& Data Engineering :} Apache Airflow (Directed Acyclic Graphs -- DAGs), Apache Kafka (Streaming temps réel), Architecture Medallion Lakehouse (Bronze/Silver/Gold), Pipelines ETL/ELT (Extract, Transform, Load), PostgreSQL, Neo4j, MySQL, Apache Hive.
  \item \textbf{Développement Logiciel \& Cloud / DevOps :} Python, C/C++, Java, JavaScript, Frameworks FastAPI et Node.js/Express, Conteneurisation Docker, Docker Compose, Intégration et Déploiement Continus (CI/CD), Flutter, React, React Native.
\end{itemize}

%=============================================================================
% 3. EXPÉRIENCES PROFESSIONNELLES (FORMULE GOOGLE X-Y-Z & RÉSULTATS CHIFFRÉS)
%=============================================================================
\cvsection{\faBriefcase}{Expériences Professionnelles}

\textbf{\color{primary} Stagiaire Ingénieur -- Data \& Systèmes Industriels} \hfill {\scriptsize \color{textmuted} \safeicon{\faCalendar*\ }07/2026 -- 09/2026}\\
{\footnotesize \textbf{\color{accent} COSUMAR (Compagnie Sucrière Marocaine)}} \hfill {\scriptsize \color{textmuted} \safeicon{\faMapMarker*\ }Casablanca, Maroc}
\begin{itemize}
  \item Conception et automatisation de pipelines ETL (Extract, Transform, Load) en Python et SQL, réduisant le temps de génération des rapports d'activité de 75\% pour plus de 20 indicateurs clés de performance (KPIs -- Key Performance Indicators).
  \item Centralisation et traitement des flux de données de production et maintenance industrielle, fiabilisant le suivi opérationnel en temps réel.
  \item Élaboration de tableaux de bord décisionnels interactifs, permettant aux équipes d'exploitation de réduire de 30\% le temps de diagnostic des anomalies.
  \item \textit{Technologies :} Python, SQL, Pipelines ETL, Tableaux de bord décisionnels, Systèmes industriels.
\end{itemize}

\vspace{0.5pt}

\textbf{\color{primary} Stagiaire Développeur Solutions Numériques} \hfill {\scriptsize \color{textmuted} \safeicon{\faCalendar*\ }07/2025 -- 09/2025}\\
{\footnotesize \textbf{\color{accent} SEWS-CABINET (Sumitomo Electric Group)}} \hfill {\scriptsize \color{textmuted} \safeicon{\faMapMarker*\ }Casablanca, Maroc}
\begin{itemize}
  \item Développement et déploiement de l'application industrielle collaborative SEWS Connect sous Flutter et Riverpod pour plus de 50 opérateurs d'atelier.
  \item Intégration d'un module d'appels vidéo et audio temps réel (WebRTC \& Agora RTC), assurant des communications stables avec une latence $<$ 250 ms.
  \item Implémentation d'un système de scan QR Code (Quick Response Code) avec base locale résiliente offline-first (Hive) synchronisée avec Firebase, diminuant le délai d'assignation des interventions d'urgence de 40\%.
  \item \textit{Technologies :} Flutter, Riverpod, WebRTC, Agora RTC, Hive, Firebase, QR Code.
\end{itemize}

%=============================================================================
% 4. PROJETS D'INGÉNIERIE PHARES
%=============================================================================
\cvsection{\faCogs}{Projets d'Ingénierie Phares}

\textbf{\color{primary} AI Multi-Agent Orchestrator -- Framework d'Agents Autonomes \& RAG} \hfill {\scriptsize \href{https://github.com/anasmouquinee/AI_MULTI_AGENT_PROJECT}{\color{accent}\safeicon{\faGithub\ }Code Source GitHub \safeicon{\faExternalLink*}}}\\
{\scriptsize \color{textmuted} Architecture multi-agents avec LangChain, FastAPI et recherche vectorielle sémantique}
\begin{itemize}
  \item Conception d'une architecture multi-agents coordonnant la recherche, la génération de code et la validation (QA), réduisant le temps de traitement de 35\%.
  \item Intégration d'un pipeline RAG (Retrieval-Augmented Generation) avec recherche vectorielle sémantique et banc de test validant une pertinence $>$ 90\%.
\end{itemize}

\vspace{0.5pt}

\textbf{\color{primary} Enterprise Media Lakehouse \& Plateforme Big Data Temps Réel} \hfill {\scriptsize \href{https://github.com/anasmouquinee/Projet-Architecture-de-donn-es}{\color{accent}\safeicon{\faGithub\ }Code Source GitHub \safeicon{\faExternalLink*}}}\\
{\scriptsize \color{textmuted} Plateforme distribuée avec Apache Airflow, Kafka et architecture Lakehouse}
\begin{itemize}
  \item Déploiement d'une architecture Medallion Lakehouse (Bronze $\to$ Silver $\to$ Gold) traitant un flux continu de plus de 100\,000 actualités et événements par jour.
  \item Orchestration et monitoring des flux de données par DAGs Apache Airflow, streaming via Kafka et conteneurisation Docker multi-services.
\end{itemize}

\vspace{0.5pt}

\textbf{\color{primary} Film-A -- Moteur de Recommandation IA par Graphes de Connaissances} \hfill {\scriptsize \href{https://github.com/anasmouquinee/film-rec}{\color{accent}\safeicon{\faGithub\ }Code Source GitHub \safeicon{\faExternalLink*}}}\\
{\scriptsize \color{textmuted} Base graphe Neo4j, Grand Modèle de Langage (LLM Gemini AI) et Mobile React Native}
\begin{itemize}
  \item Conception d'un moteur hybride associant Neo4j et l'API Gemini pour générer des recommandations cinéphiles contextualisées et explicables en temps réel.
  \item Développement de l'interface mobile sous React Native connectée à une API REST (Application Programming Interface) sous Node.js.
\end{itemize}

%=============================================================================
% 5. FORMATION ACADÉMIQUE
%=============================================================================
\cvsection{\faGraduationCap}{Formation Académique}

\textbf{\color{primary} Cycle Ingénieur d'État en Informatique et Réseaux} \hfill {\scriptsize \color{textmuted} \safeicon{\faCalendar*\ }2022 -- 2026}\\
{\footnotesize \textbf{École Marocaine des Sciences de l'Ingénieur (EMSI)}} -- Option : \textbf{Intelligence Artificielle (IA) \& Data Science} \hfill {\scriptsize \color{textmuted} \safeicon{\faMapMarker*\ }Casablanca, Maroc}

\vspace{0.5pt}

\textbf{\color{primary} Baccalauréat Scientifique} \hfill {\scriptsize \color{textmuted} \safeicon{\faCalendar*\ }2021 -- 2022}\\
{\footnotesize \textbf{Lycée Al Anbar}} -- Sciences Physiques et Chimiques \hfill {\scriptsize \color{textmuted} \safeicon{\faMapMarker*\ }Casablanca, Maroc}

%=============================================================================
% 6. CERTIFICATIONS PROFESSIONNELLES
%=============================================================================
\cvsection{\faAward}{Certifications Professionnelles}
\begin{itemize}
  \item \textbf{Agile Project Management} -- Google (2026) : Scrum, Sprints, Backlogs et Pilotage. \href{https://www.coursera.org/account/accomplishments/verify/G2064X5FGLVB}{\color{accent}[Vérifier \safeicon{\faExternalLink*}]}
  \item \textbf{Software Engineering: Architecture} -- HKUST (2025) : UML, Design Patterns, SDLC, SQA. \href{https://www.coursera.org/account/accomplishments/verify/PYISBTHNMK4R}{\color{accent}[Vérifier \safeicon{\faExternalLink*}]}
  \item \textbf{Python Web Data \& APIs} -- Univ. of Michigan (2025) : Scraping, APIs REST, JSON/XML. \href{https://www.coursera.org/account/accomplishments/verify/3HIOWYMU2G7V}{\color{accent}[Vérifier \safeicon{\faExternalLink*}]}
  \item \textbf{POO en C++ \& Algorithmique} -- EPFL (2025) : Gestion de la mémoire et complexité. \href{https://www.coursera.org/account/accomplishments/verify/402FN9N1EN8K}{\color{accent}[Vérifier \safeicon{\faExternalLink*}]}
  \item \textbf{Advanced React \& State Architecture} -- Meta (2024) : Architecture composants, Hooks, tests. \href{https://www.coursera.org/account/accomplishments/verify/EOQZZ5AJA4PN}{\color{accent}[Vérifier \safeicon{\faExternalLink*}]}
\end{itemize}

%=============================================================================
% 7. LANGUES & CENTRES D'INTÉRÊT
%=============================================================================
\cvsection{\faLanguage}{Langues \& Centres d'Intérêt}
\begin{itemize}
  \item \textbf{Langues :} Arabe (Langue maternelle), Anglais (Professionnel / C1), Français (Courant / Professionnel C1).
  \item \textbf{Centres d'intérêt :} Veille IA \& agents autonomes, Projets Open Source (25+ dépôts GitHub), Création numérique.
\end{itemize}

\end{document}
